\documentclass[11pt,a4paper]{article}
\usepackage[T1]{fontenc}
\usepackage[utf8]{inputenc}
\usepackage{lmodern,microtype}
\usepackage[margin=27mm]{geometry}
\usepackage{amsmath,amssymb,amsthm,mathtools}
\usepackage{booktabs,array,longtable,enumitem,xcolor}
\usepackage[colorlinks=true,linkcolor=blue!45!black,citecolor=blue!45!black,urlcolor=blue!45!black]{hyperref}
\usepackage{fancyhdr}
\pagestyle{fancy}\fancyhf{}
\fancyhead[L]{\small Derivative obstructions, polytopes, and peak transfer}
\fancyhead[R]{\small T81--T85}\fancyfoot[C]{\thepage}
\setlength{\headheight}{14pt}
\setlength{\parskip}{4pt}\setlength{\parindent}{1.1em}
\setlist{itemsep=3pt,topsep=5pt}
\newtheorem{theorem}{Theorem}
\renewcommand{\thetheorem}{T\arabic{theorem}}
\setcounter{theorem}{80}
\newtheorem{lemma}{Lemma}[section]
\newtheorem{proposition}[lemma]{Proposition}
\newtheorem{corollary}[lemma]{Corollary}
\theoremstyle{definition}\newtheorem{definition}[lemma]{Definition}
\theoremstyle{remark}\newtheorem{remark}[lemma]{Remark}
\newcommand{\Z}{\mathbb Z}\newcommand{\Q}{\mathbb Q}
\newcommand{\C}{\mathbb C}\newcommand{\R}{\mathbb R}
\hypersetup{pdftitle={Derivative Obstructions, Minimal Ehrhart Counterexamples, and Peak-Transfer Bounds: T81-T85},pdfauthor={KEIO}}
\title{Derivative Obstructions,\\Minimal Ehrhart Counterexamples,\\and Peak-Transfer Bounds\\[5pt]\large Five theorems with exact prior-problem applications}
\author{KEIO}
\date{24 October 2026\\[3pt]\small Research manuscript T81--T85}
\begin{document}
\maketitle
\begin{abstract}
We give five complete arguments in three settings. For a family of Stern
polynomials, a derivative certificate proves that the monotonicity clause
of Ulas's Conjecture~4.2(1) fails for every $n\geq13$ at the fixed parameter
$k=20$; in fact eventual monotonicity fails for every even $k\geq16$.
For equivariant Ehrhart theory, projection fibers certify an unbounded
family of minimal counterexamples in dimension ten, addressing Clarke's
Question~6.6 with a specific priority qualification. A parity boundary
formula then gives a negative answer to Clarke's Question~6.7 and proves
effectiveness of every polynomial equivariant numerator for an order-two
action in dimension three.
For quantum walks, we determine the exact minimum squared-edge-weight
cost for nonzero peak transfer at prescribed distance and time. An integral
gap-group argument transfers this to adjacency and Laplacian edge bounds
for Open Problem~7.2 of Coutinho--Guo--Schmeits. The adjacency bound has
optimal cubic leading constant $1/24$ and determines the growth order of
the minimum. All arguments are unconditional; their prior ingredients and
the limits of the historical-novelty assessment are stated explicitly.
\end{abstract}

\section{Exact targets and scope}

The workflow in this manuscript is to formulate an auxiliary construction
or optimization problem, prove a quantitative statement about it, and apply
that statement to a previously published question without changing the
question's hypotheses. A negative answer is distinguished from a proof of
the proposed assertion. Supporting results are not counted as independently
posed conjectures.

\begin{center}\small
\begin{tabular}{@{}p{.07\textwidth}p{.45\textwidth}p{.39\textwidth}@{}}
\toprule
Label & Result & Prior target or role \\
\midrule
T81 & A uniform negative derivative for an infinite terminal range of Stern polynomials. & Refutes the monotonicity clause of Ulas Conjecture~4.2(1) \cite{Ulas20}.\\
T82 & Minimal ten-dimensional counterexamples with unbounded point count, volume, and $h^*_2$. & Negative answer and quantitative strengthening for Clarke Question~6.6 \cite{Clarke26}; prior-example overlap is explicit.\\
T83 & Polynomiality excludes obstructing faces in dimension three. & Negative answer to Clarke Question~6.7 \cite{Clarke26}.\\
T84 & Exact weighted energy minimum for arbitrary static graphs with nonzero peak transfer. & Auxiliary optimization used in T85; the attaining chain and its cost are prior.\\
T85 & Adjacency cubic and Laplacian quadratic edge bounds; sharp adjacency leading constant. & Answers the lower-bound request in Open Problem~7.2 \cite{CGS25}.\\
\bottomrule
\end{tabular}
\end{center}

The numbers T81--T85 label five theorems, not five independent formerly
open problems. T81 leaves the separate one-real-root assertion untouched.
T83 does not determine the least dimension of an ineffective polynomial
numerator. T85 gives lower bounds and the adjacency growth order, not the
exact minimum edge count for each finite distance.

This is a research draft checked internally by independent derivations,
not an externally refereed paper or a formal proof-assistant development.
No theorem below assumes an unproved conjecture. Historical priority is
a different question: the proofs use classical identities, reciprocity,
joins, spectral moments, and Galois arguments. In particular, the bare
negative answer to Question~6.6 can also be extracted from a known example
once its minimality is supplied. We therefore do not certify that all of
the requested historical success conditions have been met, or that a
historically unprecedented foundational principle has been discovered.

\paragraph{Ehrhart conventions.}
For an $n$-dimensional lattice polytope $P$ at height one, write
\[
 \sum_{j\geq0}\#(jP\cap L)t^j=\frac{h^*_P(t)}{(1-t)^{n+1}}.
\]
The normalized volume is $h^*_P(1)$, and
$\#(P\cap L)=n+1+[t]h^*_P(t)$. For an involution $\sigma$, let
$\chi_0,\chi_1$ be the trivial and sign characters, and let $\rho$ be the
homogeneous ambient representation. If $E_P(t)$ records the permutation
representation on the lattice points of each dilation, its equivariant
numerator is
\[
 \varphi_P(t)=\det(I-t\rho)E_P(t)
        \in R(\langle\sigma\rangle)[[t]].
\]
Its identity evaluation is $h^*_P(t)$. A coefficient is effective precisely
when both its trivial and sign multiplicities are nonnegative. If the
ambient eigenspace dimensions are $p,q$, then
$\det(I-t\rho)(\sigma)=(1-t)^p(1+t)^q$. These are the conventions in
\cite{Clarke26}.


\section{An eventual derivative obstruction for Stern polynomials}
\label{sec:ulas-obstruction}

The target is Conjecture~4.2(1) of Ulas~\cite{Ulas20}, first circulated in
2019. Write $B_0(t)=0$, $B_1(t)=1$, and
\[
 B_{2m}(t)=tB_m(t),\qquad B_{2m+1}(t)=B_m(t)+B_{m+1}(t).
\]
For integers $k\geq2$, $n\geq1$, the source defines
\[
 p_{k,n}=2^{2n+k}-3\cdot2^{n+k-1}+2^k-3,
 \qquad V_{k,n}(t)=B_{p_{k,n}}(t).
\]
The conjecture asserts that for each even $k$ there exists $c_k$ such that,
for all $n\geq c_k$, the polynomial $V_{k,n}$ has exactly one real root and
is increasing on the entire real line. We disprove its monotonicity clause
with a fixed even $k$ and an infinite terminal range of $n$. A finite
counterexample at one $n$ would not suffice to disprove this eventual
assertion.

The auxiliary problem is to certify the eventual sign of a derivative when
the polynomial degree increases but its rational expression has only a fixed
number of powers depending on the index. The following certificate makes
the quantifier over all sufficiently large indices explicit.

\begin{lemma}[A finite derivative certificate]\label{lem:derivative-certificate}
Let $H_0,H_1,H_2$ be real differentiable functions near a fixed real number
$a$ with $0<r:=|a|<1$. Suppose
\[
 F_n(t)=H_0(t)+H_1(t)t^{n-1}+H_2(t)t^{2n-2}.
\]
Then
\begin{align*}
 |F_n'(a)-H_0'(a)|
 &\leq\left(|H_1'(a)|+\frac{(n-1)|H_1(a)|}{r}\right)r^{n-1}\\
 &\quad+\left(|H_2'(a)|+\frac{2(n-1)|H_2(a)|}{r}\right)r^{2n-2}.
\end{align*}
In particular, $F_n'(a)\longrightarrow H_0'(a)$. If $H_0'(a)<0$,
then $F_n$ is not increasing on the real line for all sufficiently large
$n$.
\end{lemma}
\begin{proof}
Differentiate the displayed identity and apply the triangle inequality.
Since $0<r<1$, both $nr^{n-1}$ and $nr^{2n-2}$ tend to zero. A differentiable
increasing real function has nonnegative derivative wherever its derivative
exists, which proves the last assertion.
\end{proof}

\begin{theorem}[Failure of Ulas's eventual monotonicity conjecture]
\label{thm:ulas-counterexample}
For every integer $n\geq13$,
\[
 B_{2^{2n+20}-3\cdot2^{n+19}+2^{20}-3}'(-3/4)<0.
\]
Consequently no constant $c_{20}$ has the monotonicity property required by
Conjecture~4.2(1) of Ulas, and that original conjecture is false.
\end{theorem}

\begin{proof}
We first record the exact polynomial expression needed for the certificate.
The Stern addition identity is
\begin{equation}\label{eq:stern-addition-ulas}
 B_{m2^a\pm r}(t)=B_{2^a-r}(t)B_m(t)+B_r(t)B_{m\pm1}(t)
 \qquad(0\leq r\leq2^a),
\end{equation}
with the minus version used only when its indices are nonnegative.
It follows by induction from the defining recurrences and is also
\cite[Lemma~2.1]{Ulas20}. The geometric-sum identities give
\[
 B_{2^k-3}(t)=\frac{2t^{k-1}-t-1}{t-1},\quad
 B_{2^{n+1}-3}(t)=\frac{2t^n-t-1}{t-1},\quad
 B_{2^{n+1}-2}(t)=\frac{t(t^n-1)}{t-1}.
\]
All quotients represent polynomials; we will evaluate only at $t=-3/4$.
Writing
\[
 p_{k,n}=2^{n+k-1}(2^{n+1}-3)+(2^k-3)
\]
and then using \eqref{eq:stern-addition-ulas} on the complementary index,
we obtain
\begin{align}
 V_{k,n}(t)
 &=\left((t+1)t^{n-1}
       +\frac{(2t^{k-1}-t-1)(t^{n-1}-1)}{(t-1)^2}\right)
       \frac{2t^n-t-1}{t-1}\notag\\
 &\quad+\frac{t(2t^{k-1}-t-1)(t^n-1)}{(t-1)^2}.
 \label{eq:ulas-exact-expression}
\end{align}
This expression already occurs in the proof of
\cite[Theorem~2.3]{Ulas20}; the derivative obstruction below is the additional
step.

Put $u=t^{n-1}$. Direct collection of its three powers gives
\begin{equation}\label{eq:ulas-three-powers}
 (t-1)^3V_{k,n}(t)=C_0(t)+C_1(t)u+C_2(t)u^2,
\end{equation}
where
\begin{align*}
 C_0(t)&=(2t^{k-1}-t-1)(-t^2+2t+1),\\
 C_1(t)&=-2t^4+6t^2+4t
              +(2t^3-2t^2-6t-2)t^{k-1},\\
 C_2(t)&=2t^4-2t^3-4t^2+4t^k.
\end{align*}
Fix $k=20$, let $H_j=C_j/(t-1)^3$, and set $a=-3/4$, $r=3/4$.
Exact rational substitution gives
\[
 H_0'(a)=-\frac{110236939137}{1473173782528}<-\frac7{100}.
\]
For the other quantities in Lemma~\ref{lem:derivative-certificate}, the
exact values and simple upper bounds are
\begin{align*}
 |H_1'(a)|&=\frac{4432120404109}{11785390260224}<\frac12,&
 \frac{|H_1(a)|}{r}&=\frac{108928627487}{1683627180032}<\frac1{10},\\
 |H_2'(a)|&=\frac{164405164071}{368293445632}<\frac12,&
 \frac{2|H_2(a)|}{r}&=\frac{9957814131}{26306674688}<\frac12.
\end{align*}
Thus
\[
 |V_{20,n}'(a)-H_0'(a)|
 <E_n:=\frac{n+4}{10}\left(\frac34\right)^{n-1}
             +\frac n2\left(\frac34\right)^{2n-2}.
\]
Both summands decrease for $n\geq2$: their consecutive ratios are,
respectively,
\[
 \frac34\frac{n+5}{n+4}<1,\qquad
 \frac9{16}\frac{n+1}{n}<1.
\]
Finally, the rational inequality
\[
 E_{13}=\frac{17}{10}\left(\frac34\right)^{12}
             +\frac{13}{2}\left(\frac34\right)^{24}<\frac7{100}
\]
shows that $V_{20,n}'(a)<0$ for every $n\geq13$. This excludes every
possible choice of $c_{20}$ and proves the theorem.
\end{proof}

The mechanism can also be seen in the limiting profile on $0<x<1$:
\[
 \lim_{n\to\infty}V_{k,n}(-x)
 =\frac{(2x^{k-1}-x+1)(1-2x-x^2)}{(1+x)^3}
 \quad\text{for even }k.
\]
Its derivative at $x=3/4$ is positive for $k=20$, which is incompatible
with eventual increase in $t=-x$. The finite certificate above upgrades
this observation to the uniform threshold $n\geq13$.

In fact the failure occurs for every even $k\geq16$. The derivative of
the limiting profile at $x=3/4$ is exactly
\[
 \frac{8}{2401}\left(36-(119k+22)\left(\frac34\right)^{k-2}\right).
\]
At $k=16$ it equals $32262723/5754585088>0$, and the sequence
$(119k+22)(3/4)^{k-2}$ decreases for integer $k\geq3$. Thus, for each
even $k\geq16$, Lemma~\ref{lem:derivative-certificate} gives a terminal
range of $n$ on which $V_{k,n}'(-3/4)<0$. Only the explicit threshold
$n\geq13$ for $k=20$ is needed for the theorem.

The result refutes the entire conjunction in the stated conjecture through
one of its required clauses. It does not assert that the separate
single-real-root clause is false. The proof uses a prior Stern identity,
rational differentiation, and a geometric error estimate. These tools are
classical; this argument is not presented as a historically unprecedented
foundational principle.

\section{Projection fibers and unbounded minimal counterexamples}
\label{sec:minimal-counterexamples}

Clarke's Question~6.6 \cite{Clarke26} asks whether a minimal counterexample
to effectiveness must have $n+2$ lattice points, normalized volume four,
and zero coefficient of $t^2$ in its ordinary $h^*$-polynomial. Here
minimality means that no proper full-dimensional invariant lattice
subpolytope with the same fixed polytope remains. We answer this exact
question negatively by a family in fixed dimension ten, with all three
quantities unbounded.

The auxiliary problem is to certify such minimality from the fibers of
the averaging projection. This avoids enumerating invariant subpolytopes,
which may contain nonvertex lattice points. The underlying free-join and
pole-cancellation construction is prior: Ziegler's examples, recalled in
\cite[Example~9.2]{Clarke26}, already multiply the reflected-square series
by an ordinary $h^*$-polynomial. The contribution proposed here is the
minimality certificate and the following explicit unbounded family.

\begin{lemma}[Projection-fiber forcing]\label{lem:projection-fibers}
Let $P$ be an invariant lattice polytope with involution $\sigma$, and
let $\pi=(I+\sigma)/2$. Suppose that for each vertex $w$ of $P^\sigma$,
the lattice points of $P\cap\pi^{-1}(w)$ constitute a single
$\sigma$-orbit $O_w$, and that these orbits together contain every
vertex of $P$. Then no proper invariant lattice subpolytope $P'$ has
$(P')^\sigma=P^\sigma$.
\end{lemma}
\begin{proof}
Since $P'$ is convex and invariant, $\pi(P')=(P')^\sigma$.
A vertex $w$ of this image must be the image of at least one vertex of
$P'$. That vertex is a lattice point in $O_w$, and invariance of $P'$
forces the entire orbit $O_w$ into $P'$. Consequently $P'$ contains every
vertex of $P$, and $P'=P$.
\end{proof}

For $m\geq1$, define the five-dimensional lattice simplex
\begin{equation}\label{eq:minimal-simplex}
 Q_m=\operatorname{conv}\bigl\{0,e_1,e_2,e_3,e_4,
                   (1,-1,1,-m,2m)\bigr\}\subset\mathbb R^5.
\end{equation}
Let $S=[0,1]^2$ with involution $(x,y)\mapsto(1-x,y)$, and let
$I_m=[0,m+2]$ with trivial action. Use the standard lattice free join:
take the direct sum of the homogeneous lattices, use the sum of their
height functionals as height, and take the convex hull of the embedded
height-one polytopes. Put
\[
 R_m=S*Q_m*I_m,
\]
with the involution trivial on $Q_m$ and $I_m$.

\begin{theorem}[An unbounded negative answer to Clarke's Question 6.6]
\label{thm:unbounded-minimal}
For every $m\geq1$, the polytope $R_m$ has dimension ten and is minimal
in the exact sense of \cite[Question~6.6]{Clarke26}. Its equivariant
numerator is polynomial but not effective. More precisely, putting
\[
 q_m(t)=1+mt-t^2+m(m+1)t^3,
\]
its trivial and sign characters $\chi_0,\chi_1$ occur with the virtual
multiplicities
\begin{equation}\label{eq:minimal-character}
 \varphi_{R_m}(t)
   =(1+t+t^2)q_m(t)\chi_0+tq_m(t)\chi_1.
\end{equation}
The sign multiplicity at degree three equals $-1$, whereas
\begin{equation}\label{eq:minimal-invariants}
 \#(R_m\cap L)=m+13,\qquad
 \operatorname{Vol}_{\mathrm{norm}}(R_m)=4m(m+2),\qquad
 [t^2]h^*_{R_m}(t)=2m.
\end{equation}
Thus none of the three proposed necessary properties in Question~6.6
holds, and their failure is unbounded in fixed dimension.
\end{theorem}
\begin{proof}
The homogeneous vertex matrix of $Q_m$ has determinant of absolute
value $2m$. The barycentric vector
\[
 \frac1{2m}(m,-1,1,-1,m,1)
\]
maps to the integral homogeneous point $(e_5,1)$, with the six
coordinates ordered as in \eqref{eq:minimal-simplex}. Its last coordinate
has exact denominator $2m$, so it generates the quotient of order $2m$.
The fundamental parallelepiped lattice points are therefore represented
by the fractional parts of its multiples $k$, $0\leq k<2m$.

At $k=0$, the resulting height is zero. For nonzero even $k$, each of
the two pairs with opposite weights sums to one, while the two
$m$-weighted coordinates vanish. The height is two, for $m-1$ values of
$k$. For odd $k$, the latter two coordinates sum to one, so the height is
three, for $m$ values of $k$. The fundamental parallelepiped formula gives
\begin{equation}\label{eq:minimal-simplex-hstar}
 h^*_{Q_m}(t)=1+(m-1)t^2+mt^3=(1+t)(1-t+mt^2).
\end{equation}
In particular $Q_m$ has exactly its six vertices as lattice points.

The reflected square has
\[
 \varphi_S(t)=\frac{1+t+t^2}{1+t}\chi_0
                 +\frac{t}{1+t}\chi_1,
\]
as also recorded in \cite[Examples~2.3 and~2.8]{Clarke26}.
The interval has ordinary numerator $1+(m+1)t$. Under the standard
lattice join, cone semigroups are direct products and the ambient
representations are direct sums. Thus Hilbert series and determinant
factors multiply, and so do equivariant numerators. Applying
\eqref{eq:minimal-simplex-hstar} cancels the square denominator and gives
\eqref{eq:minimal-character}, since
\[
 (1-t+mt^2)(1+(m+1)t)=1+mt-t^2+m(m+1)t^3.
\]
The degree-three sign multiplicity is $-1$, proving failure of
effectiveness despite polynomiality.

The dimension of the three-factor free join is
$2+5+1+2=10$. Evaluating its numerator at the identity gives
\[
 h^*_{R_m}(t)=(1+t)^2q_m(t).
\]
Since $[t]h^*_{R_m}=m+2$, the standard degree-one Ehrhart identity gives
$\#(R_m\cap L)=11+m+2=m+13$. Its degree-two coefficient is $2m$, and
its value at $1$ is $4m(m+2)$. These prove
\eqref{eq:minimal-invariants}.

It remains to check minimality, including possible subpolytopes that use
nonvertex lattice points. The fixed polytope is the free join
$S^\sigma*Q_m*I_m$. Each vertex belonging to $Q_m$ or $I_m$ has a
singleton fiber under $\pi$: its square-factor homogeneous height is
zero, and the involution is trivial on the other factors. The two
vertices of the fixed square segment are $(1/2,0)$ and $(1/2,1)$.
Their fibers are the two unit horizontal edges of $S$. Each contains
exactly its two lattice endpoints, and those endpoints form a single
$\sigma$-orbit. The direct-sum height coordinates prevent any other join
factor from contributing to these fibers. These fibers account for all
vertices of $R_m$. Lemma~\ref{lem:projection-fibers} now gives minimality,
even without imposing the full-dimensionality restriction on competing
subpolytopes.
\end{proof}

For example, $m=1$ gives a ten-dimensional polytope with $14$ lattice
points, normalized volume $12$, and $h^*_2=2$. Its numerator is
\[
 (1+2t+t^2+2t^3+t^4+2t^5)\chi_0
       +(t+t^2-t^3+2t^4)\chi_1.
\]
The exact verification bundle independently enumerates the fundamental
parallelepiped and counts lattice points in $Q_m$ and $2Q_m$ for
$1\leq m\leq12$. Those counts are $6$ and $m+20$, respectively.

\paragraph{Priority boundary.}
The known thirteen-dimensional example in
\cite[Example~9.2]{Clarke26} also admits the same fiber-minimality
argument. Thus the bare negative answer can already be extracted from
that example once its minimality is supplied. The present result provides
an explicit fixed-dimensional unbounded family and a full minimality
proof; it does not claim that the reflected-square join mechanism or all
its counterexample inputs are new. This limits any claim to a historically
new methodology, even though the original yes/no question is answered.

\section{Half-lattice boundary detection and Clarke's dimension-three question}
\label{sec:half-lattice-boundary}

Clarke's Question~6.7 \cite{Clarke26} asks whether a three-dimensional
lattice polytope with an involution of integral type $(2,0,1)$ can have a
polynomial equivariant numerator while possessing an obstructing face.
We give a negative answer to this exact question. Our auxiliary problem is
to detect, by a boundary coefficient, an edge of a half-lattice polygon
which has no lattice point. This detection prevents the period collapse
that would be required by the proposed example.

The noncollapse of the second Ehrhart coefficient is established in the
literature; see Beck--Sam--Woods \cite{BeckSamWoods08}. We give the
special parity formula needed here directly. The additional geometric
step connects the source's nonsaturated face-generated groups to that
formula without replacing them by saturated lattices.

\subsection{A segment lemma and an exact parity formula}

For a lattice $\Lambda$, write $\tfrac12\Lambda$ for its half-lattice.
The same notation on an affine lattice means that an affine lattice origin
has been chosen.

\begin{lemma}[Half-lattice segment detection]\label{lem:half-segment}
Let $E=[u,v]$ be a nondegenerate segment with
$u,v\in\tfrac12\Lambda$. If its affine line contains a point of
$\Lambda$, then $E$ contains a point of $\Lambda$.
\end{lemma}
\begin{proof}
Choose $x_0\in\Lambda\cap\operatorname{aff}(E)$ and a primitive lattice
vector $w$ parallel to $E$. Since $2(u-x_0),2(v-x_0)\in\Lambda\cap\mathbb Rw
=\mathbb Zw$, there are distinct integers $a,b$ with
$u=x_0+(a/2)w$ and $v=x_0+(b/2)w$. A closed interval with distinct
half-integer endpoints contains an integer. If the endpoint parities
differ, one endpoint is integral; if both endpoints are odd halves, their
distance is at least one. The resulting integral parameter gives a lattice
point of $E$.
\end{proof}

Let $R$ be a two-dimensional polygon with half-lattice vertices in an
affine rank-two lattice $\Lambda$. Its Ehrhart constituents have the form
\[
 N_{\mathrm{even}}(k)=Ak^2+B_{\mathrm{even}}k+C_{\mathrm{even}},\qquad
 N_{\mathrm{odd}}(k)=Ak^2+B_{\mathrm{odd}}k+C_{\mathrm{odd}}.
\]
For an edge $E$, let $\lambda(E)>0$ denote its length measured in the
primitive lattice direction parallel to it. Write its supporting line as
$\langle n_E,x\rangle=c_E$, where $n_E$ is primitive integral.
The half-lattice vertices imply $2c_E\in\mathbb Z$.

\begin{lemma}[Parity boundary formula]\label{lem:parity-boundary}
With this notation,
\begin{equation}\label{eq:parity-boundary}
 B_{\mathrm{even}}-B_{\mathrm{odd}}
   =\frac12\sum_{E:\,c_E\notin\mathbb Z}\lambda(E).
\end{equation}
In particular, a non-lattice supporting line forces a strictly positive
difference between the linear Ehrhart coefficients.
\end{lemma}
\begin{proof}
Rational Ehrhart reciprocity in dimension two gives
$\#(\operatorname{relint}(kR)\cap\Lambda)=N(-k)$. Since $k$ and $-k$
have the same parity, subtraction shows that the number of boundary
lattice points is exactly $2B_{\mathrm{even}}k$ for even positive $k$,
and $2B_{\mathrm{odd}}k$ for odd positive $k$.

For even $k$, every line
$\langle n_E,x\rangle=kc_E$ contains lattice points, and the number on
$kE$ is $k\lambda(E)+O(1)$. For odd $k$, this remains true when
$c_E\in\mathbb Z$; when $c_E\notin\mathbb Z$, the supporting line
contains no lattice point. Summing over edges overcounts vertices by a
bounded amount. Comparing linear terms in the two parity classes yields
\[
 2B_{\mathrm{even}}=\sum_E\lambda(E),\qquad
 2B_{\mathrm{odd}}=\sum_{E:\,c_E\in\mathbb Z}\lambda(E).
\]
Their difference proves \eqref{eq:parity-boundary}.
\end{proof}

\subsection{The unchanged prior question}

Use Clarke's homogeneous convention: $L$ is a rank-four lattice with a
primitive height functional, $P$ is a three-dimensional lattice polytope
at height one, and $\sigma$ is a height-preserving involution.
Type $(2,0,1)$ means that its integral normal form has two trivial
one-dimensional blocks and one block exchanging two coordinates. The
ambient eigenspace dimensions are therefore $p=3$ and $q=1$.

For an invariant face $F$, let
$M_F=\mathbb Z\langle F\cap L\rangle$. The face is \emph{$\sigma$-odd}
when every element of $M_F^\sigma$ has even height, and is
\emph{obstructing} when it is $\sigma$-odd and
$\dim F^\sigma\geq q$. These are the original definitions
\cite[Definition~4.1]{Clarke26}; $F=P$ is permitted and $M_F$ need not be
saturated. Write $\varphi_P(t)$ for the equivariant numerator.

\begin{theorem}[Negative answer to Clarke's Question 6.7]
\label{thm:clarke-dimension-three}
For every three-dimensional lattice polytope with an involution of type
$(2,0,1)$, polynomiality of $\varphi_P(t)$ implies that $P$ has no
obstructing face. Consequently no polytope requested in
\cite[Question~6.7]{Clarke26} exists.

More generally, for every three-dimensional lattice polytope with an
order-two action, polynomiality of its equivariant numerator implies the
existence of an invariant linear system of parameters of the prescribed
ambient representation type, and hence effectiveness of that numerator.
\end{theorem}
\begin{proof}
Let $R=P^\sigma$. Its dimension is two and its vertices are orbit averages
$(v+\sigma(v))/2$, by \cite[Lemma~2.2]{Clarke26}. The fixed lattice
$L^\sigma$ has rank three, and its height group is either $\mathbb Z$ or
$2\mathbb Z$.

If this height group is $2\mathbb Z$, the fixed lattice-point count is
zero at every odd dilation. At the even dilations it is a degree-two
polynomial with positive leading coefficient. The fixed Ehrhart series
therefore has a genuine pole of order three at $t=-1$. The character
identity
\begin{equation}\label{eq:three-dimensional-fixed-character}
 \varphi_P(t)(\sigma)
     =(1-t)^3(1+t)\sum_{k\geq0}\#(kR\cap L)\,t^k
\end{equation}
cancels only one order of this pole, contradicting polynomiality.

We may thus suppose that the height group is $\mathbb Z$. Choose a fixed
height-one lattice point as affine origin; it need not lie in $P$.
This identifies the fixed height-one affine lattice with a rank-two
lattice $\Lambda$, and the orbit-average formula shows that $R$ has
half-lattice vertices.

Suppose that $F$ is an obstructing face. The face contains no
$\sigma$-fixed ambient lattice point: such a point would itself belong to
$F\cap L\subset M_F$ and have height one, contrary to $\sigma$-oddness.
This implication makes no saturation assumption on $M_F$.

If $F\ne P$, then $F^\sigma$ is a proper face of $R$ by
\cite[Lemma~2.2(b)]{Clarke26}. Its dimension is at least $q=1$, so it is
an edge. It contains no lattice point. Lemma~\ref{lem:half-segment} implies
that its supporting affine line also contains no lattice point, and
Lemma~\ref{lem:parity-boundary} gives
$B_{\mathrm{even}}-B_{\mathrm{odd}}>0$.
If $F=P$, then $R$ contains no lattice point at all, so the same argument
applies to each of its edges and gives the same strict inequality. This
also handles a nonsaturated group generated by all lattice points of $P$.

The leading area coefficients of the two Ehrhart constituents agree.
Their difference is thus a degree-one polynomial with positive linear
coefficient. Write the fixed count as
\[
 N(k)=\frac{N_{\mathrm{even}}(k)+N_{\mathrm{odd}}(k)}2
       +(-1)^k\frac{N_{\mathrm{even}}(k)-N_{\mathrm{odd}}(k)}2.
\]
The generating function of the first term has poles only at $t=1$.
The second term has a genuine pole of order two at $t=-1$, since its
linear coefficient is nonzero. Formula
\eqref{eq:three-dimensional-fixed-character} cancels just one order, so
$\varphi_P(t)(\sigma)$ is not a polynomial. This contradiction rules out
every obstructing face and proves the exact negative answer.

By \cite[Theorem~4.2]{Clarke26}, absence of obstructing faces gives the
claimed linear system of parameters and effectiveness. For the extension
to all order-two actions in dimension three, if $p\leq q$ the same
conclusion is \cite[Corollary~4.4]{Clarke26}. The remaining possibility is
$(p,q)=(3,1)$. Besides type $(2,0,1)$, its only integral type is $(3,1,0)$;
in that type the fixed polytope is integral, and
\cite[Corollary~4.5]{Clarke26} applies. These cases exhaust the possible
nontrivial actions. The trivial action is the ordinary Ehrhart case.
\end{proof}

The final dimension-three statement is a consequence of the one new
prior-problem application, not another independently posed question.
It raises the order-two lower bound in Clarke's Question~6.8 from three
to four, but does not determine the least counterexample dimension:
dimensions four and five are not settled here.

The finite verification accompanying this argument uses exact integer and
rational arithmetic on $52$ half-lattice polygons. It checks both Ehrhart
constituents and \eqref{eq:parity-boundary}, including a polygon with period
collapse at a nonintegral vertex. These checks support the implementation
of the formula; the preceding lemmas provide its proof for every polygon.

\section{An exact energy problem for peak state transfer}
\label{sec:peak-energy}

Coutinho, Guo, and Schmeits ask in Open Problem~7.2 of
\cite{CGS25} for a lower bound on the number of edges of a graph admitting
peak state transfer between vertices at distance $d$. The statement does
not specify adjacency or Laplacian dynamics, so both are treated below.
Our auxiliary problem first permits real edge weights: at a fixed transfer
time and distance, what is the least possible sum of squared edge weights?
It has an exact answer. A Galois argument then removes the time parameter
for the original unweighted problem.

The ingredients have established antecedents. Coutinho's size bounds
\cite{Coutinho18} use spectral spacing and trace moments for periodic
vertices. The arithmetic characterization of mutual spectral support is in
\cite[Theorem~5.4]{CGS25}. The chain attaining our weighted bound is the
classical hypercube quotient underlying engineered perfect transfer
\cite{Christandl04}; its fixed-energy transfer time, including the same
constant, is explicitly recorded in \cite[equation~(17)]{Stepanenko25}.
The proposed contribution is the exact optimization
for arbitrary nonzero peak transfer, followed by the bounds and sharpness
for the original sparsity question; these antecedents are not presented
as a new foundational principle.

For a real symmetric matrix $M$, let
$M=\sum_\theta\theta E_\theta$. For distinct indices $u,v$, put
\[
 S=S_M(u,v)=\{\theta:(E_\theta)_{uv}\ne0\},\qquad
 c_\theta=(E_\theta)_{uv},\qquad B_M(u,v)=\sum_{\theta\in S}|c_\theta|.
\]
Nonzero peak transfer at time $\tau>0$ means
\[
 \left|(e^{i\tau M})_{uv}\right|=B_M(u,v)>0.
\]
Equality in the triangle inequality is equivalent to the existence of a
unit complex number $\gamma$ such that
\begin{equation}\label{eq:peak-alignment}
 e^{i\tau\theta}=\gamma\,\operatorname{sgn}(c_\theta)
 \qquad(\theta\in S).
\end{equation}
In particular, any two supported eigenvalues differ by an integer multiple
of $\pi/\tau$.

\begin{lemma}[Distance and mutual support]\label{lem:mutual-moments}
Suppose the graph of nonzero off-diagonal entries of $M$ has distance $d$
between $u$ and $v$, and $B_M(u,v)>0$. Then $|S_M(u,v)|\geq d+1$.
\end{lemma}
\begin{proof}
For $0\leq j<d$, a product contributing to $(M^j)_{uv}$ cannot traverse
the required $d$ off-diagonal steps, even if diagonal entries are allowed.
Thus
\[
 0=(M^j)_{uv}=\sum_{\theta\in S}c_\theta\theta^j\qquad(0\leq j<d).
\]
If $s=|S|\leq d$, the first $s$ equations form a square Vandermonde
system on the distinct elements of $S$. Its invertibility would force
all $c_\theta$ to be zero, a contradiction.
\end{proof}

\begin{lemma}[Quadratic cost of separated points]\label{lem:gap-energy}
If $x_1<\cdots<x_s$ and $x_{j+1}-x_j\geq a>0$, then
\begin{equation}\label{eq:separated-energy}
 \sum_{j=1}^s x_j^2\geq\frac{a^2s(s^2-1)}{12}.
\end{equation}
Equality holds precisely for $x_j=a(j-(s+1)/2)$.
\end{lemma}
\begin{proof}
Writing $\bar x=s^{-1}\sum_jx_j$ gives
\[
 \sum_jx_j^2=s\bar x^2+
       \frac1s\sum_{i<j}(x_j-x_i)^2
 \geq\frac{a^2}{s}\sum_{i<j}(j-i)^2
 =\frac{a^2s(s^2-1)}{12}.
\]
Equality requires $\bar x=0$ and all consecutive gaps equal to $a$,
which is also sufficient.
\end{proof}

\begin{theorem}[Exact fixed-time energy minimum]\label{thm:peak-energy}
Let $G$ be a finite loopless graph with nonzero real edge weights $w_e$
that are constant in time, and let $A$ be its weighted adjacency matrix. If nonzero peak transfer
occurs across distance $d\geq1$ at time $\tau>0$, then
\begin{equation}\label{eq:peak-energy-bound}
 \tau^2\sum_{e\in E(G)}w_e^2
      \geq \frac{\pi^2d(d+1)(d+2)}{24}.
\end{equation}
For every $d\geq1$ and $\tau>0$, equality is attained by a weighted path
on $d+1$ vertices with consecutive weights
\begin{equation}\label{eq:krawtchouk-weights}
 w_j=\frac{\pi}{2\tau}\sqrt{j(d+1-j)},\qquad 1\leq j\leq d.
\end{equation}
Thus the right side is the exact optimum over all such weighted graphs,
with no assumed lower bound on the peak amplitude other than positivity.
\end{theorem}
\begin{proof}
Put $s=|S_A(u,v)|$. By Lemma~\ref{lem:mutual-moments}, $s\geq d+1$;
by \eqref{eq:peak-alignment}, its ordered eigenvalues have gaps at least
$\pi/\tau$. Each supported eigenvalue occurs in the full spectrum with
multiplicity at least one. Therefore
\[
 2\sum_ew_e^2=\operatorname{tr}(A^2)
 \geq\sum_{\theta\in S}\theta^2
 \geq\frac{\pi^2s(s^2-1)}{12\tau^2}
 \geq\frac{\pi^2d(d+1)(d+2)}{12\tau^2}.
\]
This proves the lower bound.

For attainment, put $a=\pi/\tau$ and consider on $(\mathbb C^2)^{\otimes d}$
the operator $(a/2)\sum_{j=1}^d X_j$, where $X_j$ exchanges the two
basis vectors in tensor coordinate $j$. The subspace of vectors constant
on each Hamming-weight class is invariant. In the normalized class basis,
its restriction is the path with weights \eqref{eq:krawtchouk-weights}.
At time $\tau$ each tensor factor evolves as
$e^{i\pi X/2}=iX$, so the all-zero vector is sent to $i^d$ times the
all-one vector. These two vectors are the endpoint class vectors.
Consequently the weighted path has perfect, hence peak, transfer at
$\tau$. Finally,
\[
 \sum_{j=1}^d w_j^2
   =\frac{\pi^2}{4\tau^2}\sum_{j=1}^d j(d+1-j)
   =\frac{\pi^2d(d+1)(d+2)}{24\tau^2}.
\]
The construction and its energy value are established. The lower bound
here proves optimality within the larger class of arbitrary static graphs
and nonzero peak transfer. It does not apply to time-dependent controls,
which can be faster \cite{Stepanenko25}.
\end{proof}

\section{Edge bounds for the original peak-transfer problem}
\label{sec:peak-edges}

The previous optimization still contains the time. For an integer matrix,
the mutual support supplies an arithmetic normalization even when neither
vertex is periodic. This distinction is essential: peak transfer does not
in general imply periodicity at either endpoint.

\begin{lemma}[The integral gap group]\label{lem:integral-gap-group}
Let $M$ be a symmetric integer matrix and suppose nonzero peak transfer
occurs between $u\ne v$. The additive subgroup generated by the differences
$\theta-\theta'$ with $\theta,\theta'\in S_M(u,v)$ is
$\beta\mathbb Z$, where $\beta>0$ and $\beta^2$ is a positive integer.
The positive peak times are precisely the odd multiples of $\pi/\beta$.
If additionally $S_M(u,v)\subseteq[0,\infty)$ and $0\in S_M(u,v)$,
then $\beta$ is a positive integer and all elements of $S_M(u,v)$ are
nonnegative integer multiples of $\beta$.
\end{lemma}
\begin{proof}
Let $K\subset\mathbb R$ be the splitting field of the characteristic
polynomial. Its roots are real algebraic integers. The spectral-projector
formula
\[
 E_\theta=\prod_{\lambda\ne\theta}
                 \frac{M-\lambda I}{\theta-\lambda}
\]
shows that every Galois automorphism of $K/\mathbb Q$ sends $E_\theta$
to $E_{\sigma\theta}$. Hence $S_M(u,v)$ is Galois-stable.

At an existing peak time $\tau$, all differences belong to
$(\pi/\tau)\mathbb Z$. The group they generate is consequently
$\beta\mathbb Z$ for some $\beta>0$, and its positive generator is an
integer linear combination of those differences. Thus $\beta$ is an
algebraic integer. Galois stability of $S$ makes this group invariant;
its only two generators are $\beta$ and $-\beta$. Therefore
$\sigma\beta=\pm\beta$ for every $\sigma$, and $\beta^2$ is rational
and an algebraic integer. It is a positive integer.

Fix $\theta_0\in S$ and write $\theta=\theta_0+n_\theta\beta$.
The integers $n_\theta$ have greatest common divisor one. The existing
peak alignment implies $\tau\beta=k\pi$ for an integer $k$. Since
$\sum_{\theta\in S}c_\theta=(I)_{uv}=0$ and not all $c_\theta$ vanish,
both signs occur. Thus $k$ is odd: an even $k$ would give the same phase
at every supported eigenvalue. At $\pi/\beta$ the phase ratios
$(-1)^{n_\theta}$ are the same as at $\tau$, proving peak transfer there.
The same reasoning proves exactly the asserted set of peak times.

For the last claim, $0\in S$ gives $S\subseteq\beta\mathbb Z$. If some
Galois automorphism sent $\beta$ to $-\beta$, it would send a positive
element of $S$ to a negative element, contradicting $S\subseteq[0,\infty)$.
Such a positive element exists because $u\ne v$ and $|S|\geq2$.
Therefore $\beta$ is fixed by Galois, and is a positive rational algebraic
integer, hence a positive integer.
\end{proof}

\begin{theorem}[Peak-transfer sparsity bounds and cubic optimality]
\label{thm:peak-edge-bounds}
Let a finite simple graph have $m$ edges and nonzero peak transfer between
vertices at graph distance $d\geq1$.
\begin{enumerate}
\item For adjacency dynamics,
\begin{equation}\label{eq:adjacency-edge-bound}
 m\geq\frac{d(d+1)(d+2)}{24}.
\end{equation}
The leading constant $1/24$ is optimal: an infinite sequence of unweighted
graphs and vertex pairs has $d\to\infty$ and $m/d^3\to1/24$.
\item For Laplacian dynamics,
\begin{equation}\label{eq:laplacian-edge-bound}
 m\geq\frac{d(d+1)}4.
\end{equation}
No sharpness claim for the Laplacian bound is made here.
\end{enumerate}
In particular these are unconditional lower bounds answering the formal
request of \cite[Open Problem~7.2]{CGS25}. If $m_A(d)$ denotes the minimum
number of edges for adjacency peak transfer across distance $d$, then
$m_A(d)=\Theta(d^3)$. The exact finite-distance minimum is not determined.
\end{theorem}
\begin{proof}
One may restrict to the connected component containing the two vertices;
its edge count is no larger than $m$. For adjacency dynamics,
Lemma~\ref{lem:integral-gap-group} gives a peak time $\pi/\beta\leq\pi$.
Theorem~\ref{thm:peak-energy}, with all weights one, gives the stronger
bound
\[
 m\geq\frac{\beta^2d(d+1)(d+2)}{24}
      \geq\frac{d(d+1)(d+2)}{24}.
\]

For Laplacian dynamics, $L$ is positive semidefinite and its zero
projector is $J/N$ on a connected graph of order $N$. Thus $0$ is in
mutual support. By Lemma~\ref{lem:integral-gap-group}, the supported
eigenvalues are distinct nonnegative integers. If their number is $s$,
then $s\geq d+1$ by Lemma~\ref{lem:mutual-moments}, which permits the
diagonal entries of $L$. Consequently
\[
 2m=\operatorname{tr}L\geq\sum_{\theta\in S_L(u,v)}\theta
       \geq\frac{s(s-1)}2\geq\frac{d(d+1)}2.
\]
More precisely the first gap generator improves this to
$m\geq\beta d(d+1)/4$.

We give the sharpness construction and the all-distance upper bound in
the following paragraphs, completing the theorem.
\end{proof}

\subsection{An unweighted sharpness family, verified at individual vertices}

The graphs $Y_n$ in \cite[Section~6.3]{CGS25} provide the required
upper examples. For completeness, their spectral and transfer properties
needed here are verified directly. For $n\geq2$, take $2n$ independent
vertex classes in a chain, with sizes
\[
 s_{2j-1}=n-j+1,\qquad s_{2j}=j\qquad(1\leq j\leq n),
\]
and join consecutive classes by complete bipartite graphs, with no other
edges. The vertex and edge counts are
\begin{equation}\label{eq:Yn-counts}
 |V(Y_n)|=n(n+1),\qquad
 |E(Y_n)|=\sum_{j=1}^n j(n-j+1)+\sum_{j=1}^{n-1}j(n-j)
          =\frac{n(n+1)(2n+1)}6.
\end{equation}
Let $x$ be the unique vertex in class $2$, and let $y$ be any one vertex
in class $2n$. Their distance is $2n-2$.

Let $N$ have the normalized class indicators as columns. Then $N^TN=I$,
and the full adjacency matrix is $A= NQN^T$, where $Q$ is the weighted
path with off-diagonal entries $\sqrt{s_js_{j+1}}$.
We first verify that its eigenvalues are $\pm1,\ldots,\pm n$.
In the odd/even class ordering, $Q$ has form
$\left(\begin{smallmatrix}0&C\\C^T&0\end{smallmatrix}\right)$, where
\[
 C_{jj}=\sqrt{j(n-j+1)},\qquad
 C_{j+1,j}=\sqrt{j(n-j)}.
\]
The even-coordinate block $C^TC$ has diagonal $j(2n+1-2j)$ and
$(j,j+1)$ entry $(n-j)\sqrt{j(j+1)}$. Conjugating by the diagonal
matrix with entries $(-1)^j\sqrt j$ changes it to $I+\mathcal L$, where
on functions of $j\in\{1,\ldots,n\}$,
\[
 (\mathcal Lf)(j)=(n-j)(j+1)\bigl(f(j)-f(j+1)\bigr)
 +(n-j+1)(j-1)\bigl(f(j)-f(j-1)\bigr).
\]
Boundary terms have zero coefficient. On the polynomial basis
$1,j,\ldots,j^{n-1}$, $\mathcal L$ preserves degree and sends $j^k$ to
$k(k+2)j^k$ plus lower-degree terms. Therefore $I+\mathcal L$ has
eigenvalues $(k+1)^2$, $0\leq k<n$. This proves the claimed spectrum
of $Q$ and gives
\[
 p(z):=\det(zI-Q)=\prod_{k=1}^n(z^2-k^2).
\]

The $(2,2n)$ entry of $(zI-Q)^{-1}$ is
\[
 \frac{z\sqrt n\,((n-1)!)^2}{p(z)}.
\]
Indeed, the path cofactor is the product of the weights from positions
$2$ through $2n$, times the characteristic polynomial $z$ of the single
remaining first vertex. Taking residues at $\pm k$ yields
\[
 (E^Q_{\pm k})_{2,2n}
 =\frac{\sqrt n\,((n-1)!)^2}
        {2\prod_{1\leq j\leq n,\,j\ne k}(k^2-j^2)}.
\]
The corresponding projector entry in the full graph is divided by
$\sqrt n$, since the two endpoint classes have sizes $1$ and $n$.
The additional zero projector is $I-NN^T$ and has $(x,y)$ entry zero.
Thus every nonzero mutual coefficient has sign $(-1)^{n-k}$, the same
at $k$ and $-k$. At time $\pi$, multiplication by
$e^{\pm ik\pi}=(-1)^k$ gives the common sign $(-1)^n$.
This proves nonzero peak transfer from $x$ to each individual $y$, not
merely a statement about transfer to an entire class.

Equations \eqref{eq:Yn-counts} now give
\[
 \frac{|E(Y_n)|}{\operatorname{dist}(x,y)^3}
  =\frac{n(n+1)(2n+1)}{6(2n-2)^3}\longrightarrow\frac1{24}.
\]
Together with \eqref{eq:adjacency-edge-bound}, this proves optimality of
the universal leading constant, on an infinite sequence of even distances.

\subsection{Odd distances and the order of the minimum}

Peak transfer at a common time is preserved by Cartesian products.
To see this, suppose the nonzero mutual spectral coefficients in two
factors are $c_\theta$ and $b_\lambda$, aligned at $\tau$ with phases
$\gamma_1,\gamma_2$. The product adjacency matrix is
$A\otimes I+I\otimes B$. Before grouping repeated eigenvalue sums, its
mutual coefficients are $c_\theta b_\lambda$ at $\theta+\lambda$.
For a fixed sum, alignment gives
\[
 e^{i\tau(\theta+\lambda)}c_\theta b_\lambda
       =\gamma_1\gamma_2|c_\theta b_\lambda|.
\]
All nonzero summands therefore have the same sign after removing the
common phase. Grouping cannot cause cancellation; hence the product
also attains its spectral triangle bound.

For an adjacent vertex pair in the Petersen graph, the eigenvalues
$3,1,-2$ have respective mutual projector entries
$1/10,1/6,-4/15$. These follow from $E_3=J/10$,
$E_3+E_1+E_{-2}=I$, and $3E_3+E_1-2E_{-2}=A$.
At $\tau=\pi$ their phases are $-1,-1,1$, so this pair has peak transfer.
The Cartesian product of $Y_n$ with the Petersen graph consequently has
peak transfer over distance $2n-1$ and has
\[
 10\frac{n(n+1)(2n+1)}6+15n(n+1)=O(n^3)
\]
edges. The even distances are supplied by $Y_n$ itself, and distance one
by $K_2$. This proves $m_A(d)=\Theta(d^3)$ for all positive integer
distances. The product construction is used only for the order bound;
it does not assert that the constant $1/24$ is attained along odd distances.

\paragraph{Scope and verification.}
The finite verification script checks the complete unweighted adjacency
matrices for $2\leq n\leq6$: their characteristic polynomials, geodesic
moments, all mutual projector entries, and exact phase signs. It also
checks the optimal weighted-path characteristic polynomials and energy
identities for $1\leq d\leq10$. These exact checks supplement the proofs.
The result answers the lower-bound request of Open Problem~7.2 and
establishes its optimal adjacency growth order. It does not classify
all equality graphs, determine the exact minimum for every distance,
or prove a sharp Laplacian asymptotic. No claim of a historically
unprecedented general trace or Galois principle is made.

\section{Attribution and the requested success conditions}

The original targets are explicitly numbered in their primary sources:
Ulas Conjecture~4.2(1) \cite{Ulas20}, Clarke Questions~6.6 and~6.7
\cite{Clarke26}, and Coutinho--Guo--Schmeits Open Problem~7.2
\cite{CGS25}. The proofs use the original Stern indices, the original
definition of a minimal invariant subpolytope, the original nonsaturated
face-generated group, and the original graph-distance variable.

T81 is a refutation of the whole conjunction through its monotonicity
clause. The infinite range $n\geq13$ is essential: isolated small indices
could not refute a claim about all sufficiently large indices. T82 answers
the literal yes/no question, but its historical classification requires
care. Clarke's already recorded thirteen-dimensional example also becomes
minimal under the projection-fiber lemma. Thus the unbounded family and
explicit minimality certificate, rather than the raw join counterexample,
are the proposed additional content. T83 answers the entire stated
three-dimensional existence question. T85 supplies the bounds formally
requested in Open Problem~7.2 and the optimal adjacency cubic order; it
does not identify every finite-distance extremal graph.

The methodological antecedents are substantial. The Stern expression is
already in Ulas. Noncollapse of the second Ehrhart coefficient is prior
\cite{BeckSamWoods08}. Free joins and reflected-square cancellation occur
in the Ziegler examples cited by Clarke. Spectral-spacing and trace bounds
are in Coutinho \cite{Coutinho18}; arithmetic mutual-support conditions
are in \cite{CGS25}; and the attaining Krawtchouk chain and its fixed-energy
constant are prior \cite{Christandl04,Stepanenko25}. The new statements
proposed here concern the specific derivative obstruction, the unbounded
minimal family, the face-to-parity implication, and the general peak-transfer
optimization and its graph consequences. Recombining these tools is not
evidence that the underlying general principles are historically new.

Primary-source and follow-up searches through 24 October 2026 located no
explicit earlier resolution identical to T81, T83, or the complete T85
bound-and-sharpness statement. This is a bounded search result, not proof
that no such work exists. The overlap for T82 and the known attaining
benchmark for T84 are stated rather than hidden. Accordingly, the five
theorems are not reported as five successes satisfying every historical
condition in the research brief.

\paragraph{Screened branches not promoted to successes.}
An attempted general-lattice Floquet extension left higher-order Fourier
cycles uncontrolled and did not solve the PT-symmetry conjecture.
An all-residue Motzkin Hankel analysis did not establish the continued-fraction
cycle required by the original conjecture, and its main Christoffel
mechanism overlaps Chern--Shi \cite{ChernShi26}. These branches are omitted
from the theorem count. Direct square-completion and fewnomial consequences
of Ulas's formulas are likewise not promoted as new methods.

\appendix
\section{Exact verification and reproducibility}

The accompanying archive contains four Python programs and their recorded
JSON outputs. Each uses exact integer, rational, or symbolic arithmetic.
The finite checks support the implementation of the formulas; the preceding
proofs establish the infinite quantifiers.

\begin{center}\small
\begin{tabular}{@{}p{.27\textwidth}p{.64\textwidth}@{}}
\toprule
Proof component & Executed verification \\
\midrule
Stern derivative certificate & Exact rational certificate at $-3/4$, with
the uniform error threshold $n=13$; independent recursive derivatives for
$13\leq n\leq100$. The script also records classical auxiliary identities,
which are not included in the new-theorem count.\\
Minimal Ehrhart family & $1\leq m\leq12$: fundamental parallelepiped
enumeration, independent exact counts in $Q_m$ and $2Q_m$, and the character,
point-count, volume, and degree-two formulas.\\
Half-lattice boundary & $52$ half-lattice polygons: both Ehrhart constituents
and the exact parity-boundary formula, including a nonintegral-vertex
period-collapse example.\\
Peak-transfer bounds & Full unweighted adjacency matrices for $2\leq n\leq6$:
characteristic polynomials, geodesic moments, individual-vertex projectors,
and phase signs. Optimal weighted paths for $1\leq d\leq10$.\\
\bottomrule
\end{tabular}
\end{center}

Separate internal checks rederived the Stern certificate directly from the
recursion, the Ehrhart family from its homogeneous lattice, the face-to-pole
argument, and the graph bounds from mutual moments and exact gap groups.
These are internal mathematical checks, not external peer review.
The TEX source is standalone and contains its bibliography; compiling
twice with \texttt{pdflatex} requires no separate figures or bibliography
database. The archive supplies its own run instructions and dependencies.


\begin{thebibliography}{99}
\bibitem{Ulas20}
M. Ulas, \emph{Strong arithmetic property of certain Stern polynomials},
Publicationes Mathematicae Debrecen \textbf{96}(3--4) (2020), 401--422.
\href{https://doi.org/10.5486/PMD.2020.8683}{doi:10.5486/PMD.2020.8683}.
Inspected preprint: \href{https://arxiv.org/abs/1909.10844v1}{arXiv:1909.10844v1},
24 September 2019.

\bibitem{Clarke26}
O. Clarke, \emph{The equivariant Ehrhart theory of polytopes with order-two
symmetries: Part II},
\href{https://arxiv.org/abs/2610.17988v1}{arXiv:2610.17988v1},
13 October 2026.

\bibitem{BeckSamWoods08}
M. Beck, S. V. Sam, and K. M. Woods,
\emph{Maximal periods of (Ehrhart) quasi-polynomials},
Journal of Combinatorial Theory, Series A \textbf{115}(3) (2008), 517--525.
\href{https://doi.org/10.1016/j.jcta.2007.05.009}{doi:10.1016/j.jcta.2007.05.009};
\href{https://arxiv.org/abs/math/0702242}{arXiv:math/0702242}.

\bibitem{CGS25}
G. Coutinho, K. Guo, and V. Schmeits,
\emph{Peak state transfer in continuous quantum walks},
\href{https://arxiv.org/abs/2505.11986v4}{arXiv:2505.11986v4},
9 September 2026 (first preprint 2025).

\bibitem{Coutinho18}
G. Coutinho, \emph{Quantum walks and the size of the graph},
Discrete Mathematics (2018).
\href{https://doi.org/10.1016/j.disc.2018.04.011}{doi:10.1016/j.disc.2018.04.011};
\href{https://arxiv.org/abs/1802.08734}{arXiv:1802.08734}.

\bibitem{Christandl04}
M. Christandl, N. Datta, A. Ekert, and A. J. Landahl,
\emph{Perfect state transfer in quantum spin networks},
Physical Review Letters \textbf{92} (2004), 187902.
\href{https://doi.org/10.1103/PhysRevLett.92.187902}{doi:10.1103/PhysRevLett.92.187902}.

\bibitem{Stepanenko25}
A. A. Stepanenko, K. S. Chernova, and M. A. Gorlach,
\emph{Time-optimal transfer of the quantum state in long qubit arrays},
Physical Review Letters \textbf{135} (2025), 010803.
\href{https://doi.org/10.1103/mxrq-7rl4}{doi:10.1103/mxrq-7rl4};
\href{https://arxiv.org/abs/2501.11933v1}{arXiv:2501.11933v1}.

\bibitem{ChernShi26}
S. Chern and W. Shi,
\emph{Hankel determinants of Catalan-like sequences},
\href{https://arxiv.org/abs/2608.27208v1}{arXiv:2608.27208v1} (2026).
\end{thebibliography}

\end{document}
